\section{Problem Statement}
Although existing sign language recognition approaches report encouraging results on benchmark datasets, many systems remain difficult to deploy in real-time communication settings. A common limitation is the imbalance between spatial and temporal modeling: models that focus heavily on frame-level appearance may miss motion continuity, whereas sequence-heavy models may lose fine-grained hand configuration cues. This trade-off directly affects recognition quality for signs that differ by subtle movement or short-duration transitions.

Another challenge is computational efficiency. High-capacity architectures may improve offline accuracy, but they often increase latency and memory usage, reducing practical usability for continuous camera input. For assistive communication, response delay must be low enough to preserve conversational flow. In addition, ISL-specific data diversity is still evolving, and variability across signer style, speed, and environmental conditions can reduce generalization when models are trained without robust feature selection and fusion strategies.

To make the problem operationally measurable, this thesis focuses on three guiding questions: how to preserve fine-grained spatial detail while modeling temporal dynamics, how to maintain low-latency inference under realistic video conditions, and how to improve robustness against signer and environment variability without excessive model complexity. These questions define the practical boundary between purely academic performance and deployment-oriented utility.

The core problem addressed in this thesis is therefore the design of a real-time ISL recognition system that can jointly model spatial detail and temporal dynamics while remaining computationally efficient. The system must maintain stable performance under realistic input variability and produce output that can directly support communication between ISL users and non-signers.
